\chapter{Precision, calibration, and acquisition energy}
\label{ch:access-resources}

An intervention can make a current identifiable while leaving its numerical
value too uncertain to be useful. Conversely, a very precise instrument can
remain blind to the desired distinction. Experimental design must address both
questions: which observations determine the target, and which attainable error
bounds make the answer useful?

This chapter connects those questions for conserved currents. It first treats
shared calibration errors, then optimizes a declared relation between sensor
precision and acquisition energy. A sharp change occurs when irreducible
calibration errors are admitted: a design problem solvable by one minimum cut
becomes NP-complete in general. We nevertheless obtain inexpensive certificates
that bound the optimum and sometimes determine it exactly.

The ingredients---linear-programming duality, max-flow/min-cut, H\"older's
inequality, and supporting hyperplanes---are established mathematics. Their
role here is to provide a connected, checkable experimental-design method.
The energy results are conditional on an explicit sensor contract. They are
not universal thermodynamic lower bounds.

\section{Recovering a target in the presence of shared nuisance}

Let $j\in\mathbb R^m$ denote a vector of currents and let $Bj=0$ encode known
conservation laws. For a graph incidence matrix, a column has $-1$ at the tail
and $+1$ at the head of its oriented edge. Negative currents are permitted;
the orientation only fixes the sign convention. Suppose the available
observations satisfy
\begin{equation}
 y=Hj+A\eta+e,\qquad |e_i|\leq\varepsilon_i,
 \label{eq:res-observation}
\end{equation}
where $\eta\in\mathbb R^k$ is an unknown fixed calibration nuisance. Its effect
across observations is the known matrix $A$. The scientific target is $b^Tj$.

Throughout this section, the admissible parameter space is the entire linear
space $\ker B$, the nuisance is unrestricted, and the error set is the full
coordinatewise box. These assumptions define the minimax claim. Bounded
traffic, known affinities, correlated errors, or prior parameter bounds can
change the sharp answer.

\begin{theorem}[Constrained linear minimax recovery]
\label{thm:res-nuisance}
Under \eqref{eq:res-observation}, the minimum, over all estimators, of the
worst-case absolute error is
\begin{align}
 R
 &=\min_{w,\phi}\ \sum_i\varepsilon_i|w_i|
 &&\text{subject to }H^Tw+B^T\phi=b,\quad A^Tw=0,
 \label{eq:res-primal}\\
 &=\max_{h,z}\ b^Th
 &&\text{subject to }Bh=0,\quad |Hh+Az|\leq\varepsilon.
 \label{eq:res-dual}
\end{align}
If the constraints in \eqref{eq:res-primal} have no solution, $R=+\infty$.
Otherwise the finite optima are attained, and an optimal estimator is $w^Ty$.
\end{theorem}

\begin{proof}
The identities in \eqref{eq:res-primal} give
$w^Ty-b^Tj=w^Te$, whose absolute value is at most
$\sum_i\varepsilon_i|w_i|$. For any feasible $(h,z)$ in
\eqref{eq:res-dual}, the two parameter pairs $(j,\eta)=(h,z)$ and $(-h,-z)$
produce the same observation $y=0$ when their errors are respectively
$-(Hh+Az)$ and $Hh+Az$. Their targets are $b^Th$ and $-b^Th$, so every
estimator incurs maximum error at least $|b^Th|$. Linear-programming duality
equates the two bounds. Equivalently, append the exact observations $Bj=0$
to \eqref{eq:res-observation} and apply linear optimal recovery to $(j,\eta)$.
If the primal identities are infeasible, the annihilator alternative provides
$Bh=0$, $Hh+Az=0$, and $b^Th\ne0$. Scaling this direction proves unbounded
ambiguity. Finite-dimensional linear programs attain their finite optima,
including the cases of redundant constraints and zero error radii.
\end{proof}

The certificate contains the estimator weights $(w,\phi)$ and the
least-favorable perturbation $(h,z)$. Checking linear identities, error bounds,
and equality of the two objective values establishes optimality without
rerunning the optimizer. The small exact solver supplied with this release
enumerates supports and vertices; that implementation is intended for examples,
not large-scale general linear programming.

\paragraph{A shared bias can be cancelled; an arbitrary bias cannot.}
On an oriented triangle, conservation gives $j_1=j_2=j_3=J$. Consider
\[
 y_1=J+\eta+e_1,\qquad y_2=-J+\eta+e_2.
\]
The estimator $(y_1-y_2)/2$ has exact worst-case error
$(\varepsilon_1+\varepsilon_2)/2$. For radii $1/10$ and $1/5$, the result is
$3/20$. If both responses are $J+\eta$, the current remains unidentifiable.
It also remains unidentifiable if the two offsets are independently unknown.
The useful operation is a physical response change that reverses the current
contribution while preserving the shared offset. Negating an existing recorded
number reverses both contributions and creates no additional access.

\paragraph{The lower-bound currents can have positive physical rates.}
For a connected bidirectional network, fix any positive stationary distribution
$\pi$. Given a finite circulation $j$, choose symmetric edge traffics
$t_e>|j_e|/2$. On an edge $e:u\to v$, set
\[
 q^e_{uv}=\frac{t_e+j_e/2}{\pi_u},\qquad
 q^e_{vu}=\frac{t_e-j_e/2}{\pi_v}.
\]
All rates are positive, $\pi$ is stationary by conservation, and the channel
current is $j_e$. A least-favorable pair $\pm h$ can use the same $\pi$ and
the same traffics $t_e>|h_e|/2$. Thus the circulation lower bound does not rely
on negative transition rates. This construction does not imply that every
fixed device with bounded rates can realize the entire current space.

\section{The exact uncertainty of one hidden current}

Let $e_0:u\to v$ be an edge whose current cannot be measured directly. Remove
that edge from the graph. Some remaining edges have sensor records
$y_e=j_e+e_e$, with $|e_e|\leq\varepsilon_e$; other edges are unobserved.
The errors jointly range over the full box. No probabilistic independence
assumption is involved.

\begin{proposition}[Current uncertainty as a minimum cut]
\label{prop:res-current-cut}
Over the full circulation space, the exact minimax absolute error for $j_{e_0}$
is the minimum $u$--$v$ cut capacity in the graph without $e_0$, assigning
capacity $\varepsilon_e$ to a measured edge and infinite capacity to an
unobserved edge. An unobserved $u$--$v$ path makes the error unbounded. If
$e_0$ is a bridge, its stationary current is zero and needs no measurement.
\end{proposition}

\begin{proof}
Conservation expresses the target current as a signed sum of the other currents
crossing any separating cut. A cut containing only measured edges therefore
gives an estimator with error at most the sum of its radii. A maximum flow from
$v$ to $u$ through the remaining graph, closed by a flow through $e_0$, gives a
circulation $h$ with $|h_e|\leq\varepsilon_e$ on every measured edge. The pair
$\pm h$ has a common possible record and target separation $2h_{e_0}$.
Max-flow/min-cut equates these bounds. An entirely unobserved return path
supports an arbitrarily scaled invisible circulation. A bridge cannot carry
a nonzero conserved current.
\end{proof}

The same minimum-cut algorithm can optimize installation cost, error, or a
transformed parameter for acquisition energy. These are different inputs and
different scientific claims. An error capacity must not be interpreted as an
energy cost merely because the same graph algorithm processes it.

\section{A declared relation between precision and acquisition energy}

Assume that an available sensor has the calibrated error contract
\begin{equation}
 \varepsilon_e(E_e)
 =\delta_e+\sqrt{\frac{a_e}{E_e}}
 =\delta_e+\sqrt{\frac{q_e^3}{E_e}},
 \qquad q_e=a_e^{1/3}>0,\quad \delta_e\geq0,
 \label{eq:res-energy-law}
\end{equation}
where $E_e>0$ is acquisition energy in consistently declared units. The floor
$\delta_e$ represents an error that additional acquisition energy cannot
remove under this contract. Zero allocation means no reading and hence
unbounded uncertainty.

Equation \eqref{eq:res-energy-law} is an input to the design problem. It must
be justified for the estimated \emph{current}, including the response model
that converts a detector output into that current. A square-root uncertainty
law for raw photon or transition counts does not by itself establish the same
law after an uncertain dynamical inversion. Finite pulses, poorly known gains,
and systematic response errors can produce different exponents or additional
floors.

The optimization below concerns fixed sensor selections and nonadaptive,
continuous energy allocations. Integer record counts, power limits, minimum
pulse durations, and history-dependent acquisition require additional
constraints. Preparation, control, reset, calibration, installation, and
thermal-support energy are additional unless explicitly included in $E_e$.

If \eqref{eq:res-energy-law} describes a simultaneous confidence envelope
valid with probability at least $1-\alpha$, the resulting target-error bound
has that same confidence. Marginal guarantees for individual sensors do not
automatically give a joint guarantee. If the envelope is conservative, the
optimum below is the optimum certified by the declared adversarial box; it
need not be the minimum energy of every physically conceivable experiment.

\begin{theorem}[Optimal allocation on a cut]
\label{thm:res-cut-allocation}
For a nonempty sensor cut $C$, define
\[
 Q_C=\sum_{e\in C}q_e,\qquad
 \Delta_C=\sum_{e\in C}\delta_e.
\]
The minimum acquisition energy giving target-current error at most $\tau>0$
through that cut is
\begin{equation}
 E_C^*=\frac{Q_C^3}{(\tau-\Delta_C)^2},\qquad \tau>\Delta_C.
 \label{eq:res-cut-energy}
\end{equation}
It is attained by
\begin{equation}
 E_e=\frac{Q_C^2q_e}{(\tau-\Delta_C)^2},\qquad
 \varepsilon_e=\delta_e+
       \frac{(\tau-\Delta_C)q_e}{Q_C}.
 \label{eq:res-allocation}
\end{equation}
No finite energy suffices if $\tau\leq\Delta_C$. An empty cut has zero cost.
\end{theorem}

\begin{proof}
Put $x_e=\sqrt{q_e^3/E_e}>0$. Feasibility requires
$\sum_{e\in C}x_e\leq\tau-\Delta_C$. H\"older's inequality gives
\[
 Q_C=\sum_{e\in C}
       \left(\frac{q_e^3}{x_e^2}\right)^{1/3}x_e^{2/3}
 \leq\left(\sum_{e\in C}E_e\right)^{1/3}
       \left(\sum_{e\in C}x_e\right)^{2/3}.
\]
This proves the lower bound in \eqref{eq:res-cut-energy}, and
\eqref{eq:res-allocation} attains equality.
\end{proof}

\begin{corollary}[Global sensor and energy design]
\label{cor:res-global-energy}
Under the full circulation and box-error model, the optimum over fixed
allocations, arbitrary sensor subsets, and all estimators is
\begin{equation}
 E^*=\min_C\frac{Q_C^3}{(\tau-\Delta_C)^2},
 \label{eq:res-global-energy}
\end{equation}
where the minimum ranges over sensor-only separating cuts with
$\Delta_C<\tau$, with the bridge exception assigned zero energy.
\end{corollary}

\begin{proof}
Every successful allocation has some separating cut whose summed error is
at most $\tau$, by Proposition~\ref{prop:res-current-cut}. Its total acquisition
energy is at least the optimal allocation on that cut. Conversely, a cut
estimator with \eqref{eq:res-allocation} attains the displayed value.
\end{proof}

The formula also supplies a finite-energy feasibility test. Let $\rho$ be the
minimum cut sum of $\delta_e$, with unavailable sensors assigned infinite
capacity. For a nonbridge target with strictly positive $a_e$, finite energy
is possible exactly when $\tau>\rho$. At $\tau=\rho$, more acquisition can
approach the floor but cannot attain it at finite energy.

\begin{remark}[Other calibrated precision exponents]
The allocation argument is not specific to a square-root law. For a declared
contract $\varepsilon_e=\delta_e+(a_e/E_e)^{1/p}$ with $p>0$, put
$q_e=a_e^{1/(p+1)}$. The same H\"older argument gives
$E_C^*=Q_C^{p+1}/(\tau-\Delta_C)^p$ and
$E_e=Q_C^p q_e/(\tau-\Delta_C)^p$.
This algebraic extension does not establish which exponent a physical sensor
obeys. The exact allocation software in this release implements $p=2$.
\end{remark}

\section{An efficient exact case and a complexity boundary}

\subsection{Zero calibration floors}
If every $\delta_e=0$, minimizing \eqref{eq:res-global-energy} amounts to
minimizing $Q_C$. Thus one minimum cut with edge capacities $q_e$ gives
\begin{equation}
 E^*=\frac{\left(\min_C\sum_{e\in C}q_e\right)^3}{\tau^2}.
 \label{eq:res-ideal-energy}
\end{equation}
The cut chooses the sensors and \eqref{eq:res-allocation} chooses their
precision. A matching flow certifies that no alternative sensor set can have
less energy under the contract.

For rational $q_e$ and $\tau$, all quantities in this certificate are rational.
Equivalently, the input sensor family has $a_e=q_e^3$. General cube roots require
algebraic or certified interval arithmetic; the exact implementation does not
silently round them. The checker verifies the graph, capacity bounds,
conservation, cut--flow equality, and the allocation identities, against the
externally supplied sensor parameters.

\subsection{Irreducible calibration floors}
Consider the following decision problem. Given a graph, target edge, available
sensors with rational $q_e>0$ and $\delta_e\geq0$, rational accuracy $\tau>0$,
and rational acquisition budget $L\geq0$, does a fixed allocation achieve
worst-case target error at most $\tau$ with total energy at most $L$?

\begin{theorem}[Floor-aware design complexity]
\label{thm:res-floor-hardness}
For binary-encoded rational inputs, this decision problem is NP-complete.
NP-hardness holds even when all nontarget sensors are available, all calibration floors
are positive, every cut is feasible at finite energy, and the graph is a target
edge plus parallel two-edge paths with distinct internal vertices.
\end{theorem}

\begin{proof}
A separating cut is a polynomial-size witness. Compute its rational $Q_C$ and
$\Delta_C$ and check
\[
 \tau>\Delta_C,\qquad Q_C^3\leq L(\tau-\Delta_C)^2.
\]
The corresponding allocation has polynomial encoding length. An empty cut is
handled separately. Thus the decision problem is in NP.

For hardness, reduce PARTITION on positive integers $v_1,\ldots,v_n$ whose sum
is $V$. Set $\beta=V$. For each $i$, introduce the path $s$--$i$--$t$, with
sensor parameters
\[
 (q,\delta)=(\beta,\beta+v_i)
 \quad\text{on }s\text{--}i,\qquad
 (q,\delta)=(\beta+v_i,\beta)
 \quad\text{on }i\text{--}t.
\]
The target is the direct edge $s$--$t$. Every separating cut in the graph
without the target chooses one edge from each path. If $S$ indexes its
second-edge choices, then
\[
 Q=n\beta+\sum_{i\in S}v_i,\qquad
 \Delta=n\beta+V-\sum_{i\in S}v_i,
 \qquad Q+\Delta=K=2n\beta+V.
\]
Set
\[
 Q_*=n\beta+\frac V2,\qquad
 \tau=K-\frac{Q_*}{3},\qquad
 L=\frac{9Q_*}{4}.
\]
All numbers have polynomial bit length. Every cut has
$\tau-\Delta=Q-Q_*/3>0$, because $Q\geq nV>Q_*/3$. Its energy is
\[
 f(Q)=\frac{Q^3}{(Q-Q_*/3)^2}.
\]
The exact factorization
\[
 4x^3-9Q_*(x-Q_*/3)^2=(x-Q_*)^2(4x-Q_*)
\]
shows that $f(x)\geq9Q_*/4$ for $x>Q_*/3$, with equality exactly at $x=Q_*$.
There is therefore a feasible cut with energy at most $L$ precisely when
$\sum_{i\in S}v_i=V/2$. This completes the reduction.
\end{proof}

The reduction establishes NP-hardness through PARTITION; it does not establish
strong NP-hardness or rule out useful approximation, pseudopolynomial methods,
or tractable subclasses. It separates an efficiently checkable design from
an efficient general search for the optimum. The result also does not mean
that scientific access is uncomputable.

Budgeted terminal-cut problems already have a substantial complexity literature.
For example, Aissi and Mahjoub's treatment of minimum $s$--$t$ cuts with budget
constraints discusses related hardness results and earlier work.\footnote{Primary
manuscript: \url{https://www.lamsade.dauphine.fr/~mahjoub/Papers/Paper-Aissi-Mahjoub-2023.pdf}.}
The direct reduction above addresses the particular nonlinear
precision--energy objective used here; it is not a claim to have discovered
the hardness of multiobjective cuts.

\section{Global certificates when exact design is difficult}

Computational hardness does not prevent a useful guarantee on a proposed
design. For any chosen rational support parameter $\lambda>0$, assign capacities
\begin{equation}
 c_e(\lambda)=3\lambda^2q_e+2\lambda^3\delta_e,
 \label{eq:res-support-capacity}
\end{equation}
and let $W(\lambda)$ be their minimum cut value.

\begin{proposition}[Supporting-cut energy bracket]
\label{prop:res-support}
The optimum floor-aware acquisition energy obeys
\begin{equation}
 E^*\geq\max\{0,W(\lambda)-2\tau\lambda^3\}.
 \label{eq:res-support-bound}
\end{equation}
If a minimum cut returned for \eqref{eq:res-support-capacity} satisfies
$\Delta_C<\tau$, its allocation in \eqref{eq:res-allocation} provides a feasible
upper bound. Matching bounds certify global optimality.
\end{proposition}

\begin{proof}
For a feasible cut, write $s=\tau-\Delta_C>0$ and $x=Q_C/s$. Then
\[
 \frac{Q_C^3}{s^2}
 -\left(3\lambda^2Q_C-2\lambda^3s\right)
 =s(x-\lambda)^2(x+2\lambda)\geq0.
\]
Minimizing the supporting expression $3\lambda^2Q_C-2\lambda^3s$ over all cuts gives
$W(\lambda)-2\tau\lambda^3$. Allowing infeasible cuts in that minimum can
only lower it, so it remains a valid lower bound. Energy is also nonnegative.
The upper bound follows from the explicit feasible allocation.
\end{proof}

One max-flow computation supplies the lower-bound certificate. Its flow and
cut can be checked without rerunning the optimization. Repeating this procedure
for a declared finite set of support parameters gives the maximum certified
lower bound and the minimum feasible upper bound encountered. This is a
supporting-hyperplane bound, not a general exact solution of the NP-complete
problem.

The hardness construction also supplies a useful failure control. At
$\lambda=3/2$, the weights of $q_e$ and $\delta_e$ coincide. The lower bound is
exactly the PARTITION threshold $L$ for every instance. If the integers have no
equal partition, the true optimum is strictly larger. Thus a valid global
lower bound need not be attainable by any physical sensor cut.

\section{A complete design example}

Consider three vertices $s,a,t$. Besides the inaccessible target $s$--$t$,
there is one edge from $s$ to $a$ and two parallel physical channels from $a$
to $t$. Let
\[
 q_{\rm single}=4,\qquad q_{\rm pair,1}=1,\qquad
 q_{\rm pair,2}=2,\qquad \tau=\frac1{10}.
\]
The two possible sensor cuts are the single first edge or the pair of parallel
channels. With zero floors, the single sensor costs $6400$ acquisition units.
The pair costs $2700$ when its energies are $900$ and $1800$, giving error
radii $1/30$ and $1/15$. Requiring equal error radii $1/20$ instead costs $3600$.

Now give the paired sensors floors $3/50$ and $1/50$, retaining a zero floor
for the single sensor. The pair requires $67500$ units to attain the same target
accuracy. The optimum switches to the single sensor at $6400$.

\begin{center}
\begin{tabular}{lr}
\toprule
Design at target error $1/10$ & Acquisition energy\\
\midrule
Zero floors: optimal sensor and precision choice & $2700$\\
Zero floors: equal precision on the optimal pair & $3600$\\
Zero floors: fewest-sensor choice & $6400$\\
Nonzero paired floors: retaining the former pair & $67500$\\
Nonzero paired floors: revised optimal choice & $6400$\\
\bottomrule
\end{tabular}
\end{center}

At support parameter $\lambda=40$, one supporting cut certifies matching lower
and upper values of $6400$ for the floor-aware problem. This example illustrates
why measurement count, precision allocation, and calibration robustness are
different design objectives. Its values are exact results for a constructed
model, not measurements of an experimental device.

\subsection{Executable example and independent checking}
From the release's \texttt{code} directory, the following example constructs
the ideal design and checks its certificate:
\begin{verbatim}
from constructive_access.completion.graph_access import Edge, Graph
from constructive_access.completion.robust_completion import (
    minimum_ideal_energy, verify_ideal_energy,
    floor_energy_bound, verify_floor_energy_bound,
)

g = Graph(("s", "a", "t"), (
    Edge("target", "s", "t"),
    Edge("single", "s", "a"),
    Edge("pair_1", "a", "t"),
    Edge("pair_2", "a", "t"),
))
q = {"single": 4, "pair_1": 1, "pair_2": 2}
design = minimum_ideal_energy(g, "target", q, "1/10")
checked = verify_ideal_energy(g, "target", q, "1/10", design)
assert checked["total_energy"] == "2700"

floor = {"single": 0, "pair_1": "3/50", "pair_2": "1/50"}
bracket = floor_energy_bound(g, "target", q, floor, "1/10", 40)
checked = verify_floor_energy_bound(
    g, "target", q, floor, "1/10", bracket
)
assert checked["lower_energy"] == "6400"
assert checked["upper_energy"] == "6400"
\end{verbatim}

The generic nuisance interface provides two separate calls:
\begin{verbatim}
certificate = robust_target_certificate(model)
checked = verify_robust_certificate(model, certificate)
\end{verbatim}
The model fields specify conservation, observations, nuisance coefficients,
the target, and error radii. The checker tests the primal and dual identities
against that supplied model. Exact inputs are integers,
rational strings, or Python \texttt{Fraction} values; floating-point inputs
are rejected.

The test suite includes shared and independent bias, inaccessible targets,
bridge targets, zero error, corrupted certificates, calibration-induced sensor
choice reversal, PARTITION instances, comparisons with exhaustive cut search,
and strict lower-bound gaps. Its purpose is to test both successful design
and the conditions under which a conclusion must be withheld. Physical use
additionally requires validation of the observation and energy contracts.
